\documentclass[12pt, a4paper]{article}
\usepackage[utf8]{inputenc}
\usepackage{amsmath, amssymb}
\usepackage{graphicx}
\usepackage{booktabs}
\usepackage{natbib}
\usepackage{hyperref}

\title{Estimating Extreme Daily Market Losses on the Johannesburg Stock Exchange using Extreme Value Theory}
\author{Your Name}
\date{\today}

\begin{document}

\maketitle

\begin{abstract}
This study applies Extreme Value Theory (EVT) to estimate extreme one-day losses on the Johannesburg Stock Exchange (JSE). We compare Block Maxima (GEV) and Peaks-Over-Threshold (GPD) approaches against empirical and Gaussian models, assessing 99\% and 99.5\% Value at Risk (VaR). 
\end{abstract}

\section{Introduction}
% Context of JSE, the importance of extreme risk modeling, and limitations of normal distributions.

\section{Data and Methodology}
\subsection{Data Description}
% Discussion of the JSE Top40 dataset, calculation of log returns $r_t$ and losses $L_t$.

\subsection{Extreme Value Theory Framework}
% Mathematical definitions of GEV (Block Maxima) and GPD (Peaks-Over-Threshold), including the role of the shape parameter $\xi$.

\section{Exploratory Data Analysis}
% Reference to return/loss time series, QQ-plots, and empirical CDFs.
% \begin{figure}[h] \centering \includegraphics{...} \caption{...} \end{figure}

\section{Model Fitting and Diagnostics}
\subsection{Block Maxima (GEV)}
\subsection{Peaks-Over-Threshold (GPD)}
% Discussion of threshold selection using mean-excess plots.

\section{Value at Risk (VaR) Estimation and Backtesting}
% Comparison table of Empirical, Gaussian, and EVT VaR estimates.
% \begin{table}[h] \centering \begin{tabular}{...} \end{tabular} \end{table}

\section{Discussion and Limitations}
% Addressing IID assumptions, volatility clustering, and small-sample tail uncertainty.

\section{Conclusion}

\bibliographystyle{plainnat}
\bibliography{references}

\end{document}